% !TeX root = ../bnuthesis-example.tex

\chapter{绪论}

\section{示例的目的与依据}

本示例展示《北京师范大学研究生学位论文编写规则（2026 版）》要求的论文结构和常见排版元素。实际论文应围绕具体研究问题展开，说明研究背景和意义，在梳理已有工作的基础上提出研究目的、流程与方法，并在结论中总结有证据支持的结果。

本示例的说明文字、示意图片和公式用于演示格式，不构成作者已经完成的科学研究。使用时应根据真实研究内容重写各部分，并核对培养单位制定的补充规则。

\section{论文组成与内容要求}

前置部分包括封面、题名页、原创性声明与使用授权书、中英文摘要和目录。图表较多时，应在目录之后分别列出插图清单、附表清单；使用大量符号、缩略词和专门术语时，应编写主要符号对照表。

正文从绪论开始，以结论或讨论结束，各章另起页。文后依次安排参考文献、附录、索引、攻读学位期间取得的学位论文相关成果、致谢和答辩委员会名单。附录与索引为可选部分，本例保留两份附录以演示编号。

中文学位论文应使用规范汉字，术语、符号和代号在全文保持一致。专业术语和缩略词首次出现时需要解释，外文术语应在译文后括注原文全称。计量单位采用国家法定计量单位，图像和照片应完整、清晰。

\subsection{题名与摘要}

题名应准确、简明地反映论文的重要内容，一般不超过 25 个汉字。摘要应能够独立说明研究目的、方法、结果和结论，突出创新点，中英文内容对应。中文博士论文摘要一般约 2000 字，硕士论文一般约 1000 字；关键词各选 3--8 个，以逗号分隔。

\subsection{学术规范与材料核对}

论文应区分作者的研究与引用资料，全部被引文献列入参考文献表，未引用的阅读材料可另列为附录中的参考书目。学术成果、资助信息和答辩委员会名单应与真实记录或正式材料一致，模板中的占位信息必须在提交前替换。
\footnote{本页脚注用于演示上标标注和页底注释。正式论文中的脚注应提供与正文相关的必要说明。}

\section{格式约定与保留的例外}

规范要求正文采用宋体小四号、固定 20 磅行距，英文使用 Times New Roman。本例通过共享配置使用 Windows 中的宋体、黑体及 Times New Roman，数学字体保留 New Computer Modern；终稿须核对实际字体。章、节、条标题采用阿拉伯数字层级，图和表分别连续编号，附录中的图表公式另行编号。编写规则前文允许图表分章编号，后文格式表要求连续编号；本例采用后文较具体的格式表。

本例按已确认的使用习惯保留校徽页、书脊插页及原有声明调用，并继续以著者排序的文献表配合正文数字上标引用。这些保留项应在正式提交前结合培养单位及印制要求核对。A3 装订封皮及其书脊属于独立印制内容，不能仅将 A4 页面放大代替。
\footnote{规则要求每章另起页，但没有要求从奇数页开始；本例不据此额外插入章间空白页。}
